我目前是浙江理工大学物理学专业硕士研究生，师从宋昌盛老师，预计于2027年6月毕业。硕士阶段，我围绕低维功能材料、多铁耦合与相关量子物性开展第一性原理及多尺度模拟研究。持续的科研训练使我逐渐形成一套明确的工作方式：从具体物理问题出发选择方法，建立可复现的计算流程，再通过交叉验证把结构、电子态、微观相互作用与宏观功能性质联系起来。我希望在博士阶段进一步研究铁电与反铁电材料、外延异质结构和极化调控机制，并逐步建立理论计算、数据方法与实验表征相互反馈的研究路径。

\section{科研训练与代表性工作}

硕士阶段，我逐步掌握了VASP、Wannier90/WannierTools、TB2J和Spirit等研究工具，能够独立开展结构优化、电子结构、自旋轨道耦合、铁电极化与翻转路径、磁交换作用、磁各向异性、Dzyaloshinskii--Moriya相互作用、Berry曲率及拓扑性质等计算，并根据研究问题构建有效模型和开展原子自旋模拟。

在二维多铁材料TcIrGe$_2$Se$_6$的研究中，我从结构稳定性和铁电翻转出发，依次分析磁基态、谷自由度、Berry曲率及微观交换机制，并利用Wannier紧束缚模型追踪Tc--Se--Ir--Se--Tc长程超超交换通道，把局域结构畸变、轨道杂化与磁交换联系起来。该工作形成了“结构—极化—电子结构—微观相互作用—有效模型—功能响应”的完整研究链条，相关论文已被\textit{Physical Review B}接收。

在另一项材料设计工作中，我研究重元素替位对磁基态、铁电稳定性、轨道杂化和交换机制的影响，并将第一性原理获得的微观参数用于有限温度模型分析。目前，我还在推进二维铁电/磁性异质结构研究，重点关注极化翻转、界面构型、层间耦合与应变对功能性质的影响。这些经历使我的研究兴趣从单一材料性质预测，逐步延伸到亚稳态、界面结构和真实可实现性等更具挑战的问题。

\section{方法创新与编程能力}

我重视在现有方法不足时主动开发分析工具。在研究TcIrGe$_2$Se$_6$长程交换机制时，常规能带和投影态密度难以定量识别具体超超交换路径。为此，我基于Wannier模型编写Python程序，解析\texttt{wannier90\_hr.dat}、Wannier中心和结构信息，并对不同原子、轨道及晶格平移之间的跃迁矩阵元进行筛选与统计，从而获得特定$d$--$p$--$d$路径的轨道杂化强度。这个过程使原本定性的机制判断转化为可计算、可比较和可复核的证据。

我也在探索机器学习与物理模型的结合。在交错磁拓扑研究中，我将Sobol参数采样、Chern数计算、分层机器学习、机制感知特征和局部相界搜索结合起来，用代理模型缩小高维参数空间中的候选区域，再返回严格的物理计算进行验证。为保证可重复性，我将项目按物理模块整理为五个代码单元，共归档34个Python源文件，并区分训练样本、外部验证样本、代表点和仍需验证的材料映射。我由此认识到，机器学习的价值在于帮助研究者更高效地发现规律并分配计算资源，而不是替代物理判断。

除Python外，我能够使用Shell/Linux完成高性能计算任务管理、参数扫描、数据提取、统计分析和可视化，并习惯记录计算参数、图件来源、报错处理与后续验证计划。程序化、模块化和文档化的工作方式，使我能够在多个长期课题并行推进时保持证据链清晰，并在多轮修改后快速恢复研究上下文。

\section{科研态度与工作方式}

我对结果可靠性和证据链完整性保持较高要求。针对U值、SOC设置、$k$点密度、Wannier拟合窗口、有限尺寸模型及边界终止方式等可能影响结论的因素，我会主动设计收敛测试和交叉验证。面对不同方法给出不一致结果、模型存在多个亚稳态或某个趋势偏离预期的情况，我通常会重新检查输入条件、结构变化和物理假设，必要时修改脚本或重建模型。我享受“发现异常—提出假设—设计对照—重新验证”的过程，并希望在博士阶段进一步提升独立提出科学问题和组织完整课题的能力。

长期科研也培养了我的项目管理与科研表达能力。我会先明确科学问题和最小验证目标，再根据计算成本安排任务顺序，将复杂课题拆解为结构、电子、极化、模型、后处理与写作等模块。目前，我已有一篇\textit{Physical Review B}论文接收，一篇相关论文在审，另有二维异质结构和机器学习方向工作持续推进。论文撰写与修改经历使我更重视图表、数据、方法和文字之间的一致性。

\section{博士阶段研究计划}

博士阶段，我希望重点研究铁电/反铁电材料、外延薄膜、异质界面和极化调控机制，关注以下相互关联的问题：极化翻转过程中的原子结构演化；铁电相、反铁电相和亚稳态之间的竞争；应变、界面、缺陷与取向对极化稳定性和翻转路径的影响；以及复杂极化畴和拓扑极化结构的形成与外场调控。

这些问题与我已有的第一性原理、NEB、界面建模、编程和机器学习基础具有连续性。同时，我希望系统学习外延薄膜制备、铁电性能测试、压电力显微镜、高分辨扫描透射电子显微镜及同步辐射结构表征等实验方法，形成“理论预测与机制假设—材料和结构设计—制备与表征—实验反馈修正模型—再设计与再验证”的研究闭环。我也希望基于第一性原理数据构建机器学习原子势，研究传统DFT难以覆盖的大尺度极化翻转、畴壁运动和有限温度结构演化，使数据方法真正服务于材料物理问题。

\section{申请动机与发展目标}

硕士阶段的科研经历让我确认，自己真正感兴趣的是发现问题、建立假设、设计验证并理解微观机制的过程。长期而言，我希望继续在高校、科研院所或高水平科研平台从事研究，围绕铁电/多铁材料、低维功能材料、界面工程与极化调控形成连续方向，并建立兼具理论计算、程序开发、机器学习和实验研究能力的个人特色。

我期待在博士阶段进入具有交叉研究氛围的团队，在扎实推进具体材料问题的同时，学习如何提炼具有普适意义的物理机制。我愿意面对新方法和新实验技术的学习成本，也愿意对异常结果持续追问。希望通过系统的博士训练，从以计算研究为主的硕士生成长为能够贯通理论预测、数据方法与实验验证的独立研究者。

% ===== 院校定制段（提交前按目标院校替换下面这段） =====
\section{院校与导师匹配（提交前定制）}

贵校在\textbf{【填写目标学科或平台优势】}方面的研究积累，以及\textbf{【填写目标导师姓名】}老师在\textbf{【填写具体研究方向】}方面的工作，与我在低维功能材料、极化调控和计算方法开发方面的基础具有良好衔接。我希望在博士阶段参与\textbf{【填写拟研究问题】}的研究，并把已有的第一性原理、多尺度模拟、编程和机器学习经验用于具体课题，同时补足实验制备与表征能力。若有机会加入该团队，我将以可验证的科学问题为中心，持续完善从模型预测到实验反馈的研究闭环。